% ============================================================================
\section{Related work}
\label{sec:related}

The taxonomies of swarm robotics are built a priori. \citet{brambilla2013swarm} classified the field by its design and
analysis methods and by collective behavior, into spatially organizing behaviors, navigation, collective
decision-making and others. \citet{schranz2020swarm} applied such classes of basic behaviors to current applications,
and \citet{rossi2018review} classified the algorithms by their mathematical structure. These taxonomies organize what a
swarm can be designed to do, not what a set of recordings contains.

Behaviors are usually compared through descriptors chosen in advance: order parameters such as polarization and
milling \citep{vicsek1995novel,couzin2002collective,dorsogna2006self}, which reveal the swarm, milling and polarized
states of fish schools \citep{tunstrom2013collective}, or the behavior characterizations of swarm robotics
\citep{gomes2013generic,gomes2014systematic}. Which groups a comparison finds depends on the pairing of these features
with a similarity measure \citep{fialho2026differ}. Classifiers learn instead human labels of flocking, alignment and
grouping \citep{abpeikar2023motion}, or the labeled phases of a real robot swarm, with which unsupervised clusterings
agree only moderately \citep{alhubaishi2026interpretable}.

Other works let behaviors emerge. Novelty search explores the controllers of limited robots, and k-medoids clusters
the archive into a taxonomy whose medoids are representative behaviors, their number chosen by inspection
\citep{brown2018discovery}. The behavior space was later learned with human feedback, and the search
extended to heterogeneous swarms and real robots \citep{mattson2023leveraging,mattson2023exploring,mattson2025discovery}.
\citet{titus2021unsupervised} compare by diffusion maps the networks that positions and velocities define over a group,
and cluster the result into three groups fixed in advance. Computational ethology goes furthest toward a taxonomy
discovered from motion, but for single animals: the postural dynamics of flies divide, by a watershed of their density,
into more than one hundred stereotyped behaviors organized hierarchically \citep{berman2014mapping,berman2016predictability},
and the approach extends to pairs of animals \citep{klibaite2017unsupervised}.

Topological data analysis summarizes shapes instead of chosen descriptors. \citet{topaz2015topological} summarized in
CROCKER plots the persistent homology of the point clouds that agents form in the space of positions and velocities. On
positions alone, such summaries classify the regimes of the D'Orsogna model better than order parameters
\citep{bhaskar2019analyzing} and separate the regimes of a crowd \citep{kern2026study}. \citet{corcoran2017modelling}
followed a school of fish by zig-zag persistence of a density estimated on a grid. Other works follow these summaries in
time \citep{xian2022capturing,sathiyakumar2024detecting}, embed persistence images \citep{suetani2025unsupervised},
learn latent dynamics of diagrams \citep{zeng2024neural}, or measure the fluctuations of a large flock across
hierarchical levels \citep{mototake2024topological}. In all of them, the topology is that of a point cloud or of a
density, and velocities, when used, are coordinates of points rather than a field. \citet{legarda2026universal}
described instead the topology of a velocity field, through the extended persistence of its Helmholtz--Hodge
potentials, and grouped the recordings into a fixed number of clusters or by HDBSCAN. \Cref{tab:related} summarizes what
each of these approaches provides.

\begin{table}[!htbp]
\centering
\caption{What each approach provides (\cmark\ yes, -- no). \emph{medoids}: k-medoids for a number of clusters chosen
by the user; \emph{partly}: with one of the clustering methods used; \emph{features}: importances of hand-crafted
features; \emph{time}: time scales; \emph{levels}: the levels of one flock.}
\label{tab:related}
\newcommand{\appr}[2]{#1\newline{\footnotesize\color{black!72}\citealp{#2}}}
\newcommand{\tw}[1]{{\footnotesize\emph{#1}}}
\newcommand{\rot}[1]{\multicolumn{1}{c}{\rotatebox[origin=l]{90}{#1}}}
\setlength{\tabcolsep}{1pt}
\renewcommand{\arraystretch}{1.0}
\newcommand{\rowsep}{\arrayrulecolor{black!22}\specialrule{0.4pt}{1.5pt}{1.5pt}\arrayrulecolor{black}}
\newcommand{\groupsep}{\arrayrulecolor{black!55}\specialrule{0.6pt}{2pt}{1.5pt}\arrayrulecolor{black}}
\newcolumntype{C}[1]{>{\centering\arraybackslash}p{#1}}
% the approach column takes what the eight narrow columns (6.81 cm) and their padding (16 x 1 pt) leave of the text
% width, so that the table is exactly as wide as the text and the shaded row is continuous
\begin{tabular}{@{}>{\raggedright\arraybackslash}p{\dimexpr\linewidth-6.81cm-16pt\relax}C{0.72cm}C{0.82cm}C{0.62cm}C{0.8cm}C{0.8cm}C{1.0cm}C{1.05cm}C{1.0cm}@{}}
\toprule
& \multicolumn{3}{c}{Representation} & \multicolumn{5}{c}{Taxonomy returned}\\
\cmidrule(lr){2-4}\cmidrule(l){5-9}
Approach & \rot{Collective} & \rot{Velocities} & \rot{Scales} & \rot{No labels} & \rot{Number} & \rot{Exemplar} &
\rot{Hierarchy} & \rot{Attribution}\\
\midrule
\multicolumn{9}{@{}l}{\emph{Taxonomies set a priori}}\\
\appr{Classes of behaviors or algorithms}{brambilla2013swarm,schranz2020swarm,rossi2018review}
  & \cmark & -- & -- & -- & -- & -- & \tw{a priori} & --\\
\groupsep
\multicolumn{9}{@{}l}{\emph{Descriptors chosen in advance}}\\
\appr{Order parameters}{vicsek1995novel,couzin2002collective,tunstrom2013collective}
  & \cmark & \cmark & -- & -- & -- & -- & -- & --\\
\rowsep
\appr{Behavior characterizations}{gomes2013generic,fialho2026differ}
  & \cmark & \tw{varies} & -- & -- & -- & -- & -- & --\\
\rowsep
\appr{Supervised classifiers and clustering}{abpeikar2023motion,alhubaishi2026interpretable}
  & \cmark & \cmark & -- & \tw{partly} & \tw{partly} & -- & -- & \tw{features}\\
\groupsep
\multicolumn{9}{@{}l}{\emph{Behaviors discovered without labels}}\\
\appr{Novelty search, then k-medoids}{brown2018discovery,mattson2023leveraging,mattson2023exploring,mattson2025discovery}
  & \cmark & \tw{varies} & -- & \cmark & -- & \tw{medoids} & -- & --\\
\rowsep
\appr{Manifold learning}{titus2021unsupervised}
  & \cmark & \cmark & -- & \cmark & -- & -- & -- & --\\
\rowsep
\appr{Behavioral maps of animals}{berman2014mapping,berman2016predictability,klibaite2017unsupervised}
  & \tw{pairs} & \cmark & \tw{time} & \cmark & \cmark & -- & \cmark & --\\
\groupsep
\multicolumn{9}{@{}l}{\emph{Topology of positions}}\\
\appr{CROCKER plots}{topaz2015topological,bhaskar2019analyzing,kern2026study}
  & \cmark & \tw{partly} & \cmark & \cmark & -- & \tw{medoids} & -- & --\\
\rowsep
\appr{Zig-zag persistence}{corcoran2017modelling}
  & \cmark & -- & -- & \cmark & -- & \tw{medoids} & -- & --\\
\rowsep
\appr{Persistence over time}{sathiyakumar2024detecting}
  & \cmark & -- & \cmark & -- & -- & -- & -- & --\\
\rowsep
\appr{Levels of a large flock}{mototake2024topological}
  & \cmark & -- & \cmark & -- & -- & -- & \tw{levels} & --\\
\groupsep
\multicolumn{9}{@{}l}{\emph{Topology of motion}}\\
\appr{Persistence of a velocity field}{legarda2026universal}
  & \cmark & \cmark & -- & \cmark & \tw{partly} & -- & -- & --\\
\groupsep
\methodrow
\textbf{This work} & \cmark & \cmark & \cmark & \cmark & \cmark & \cmark & \cmark & \cmark\\
\bottomrule
\end{tabular}
\end{table}

Each of these lines of work addresses its own problem, from organizing the design of swarms to characterizing physical
regimes, mapping the postures of single animals or comparing velocity fields, and each provides part of what a
taxonomy of collective behaviors requires. What remains open is to join a topological description of collective
motion with the unsupervised discovery of a taxonomy and with an explanation of its differences down to the agents,
which is what \cref{sec:method} presents.
